% Load transfer along a row of five fasteners. Values from a direct-stiffness model,
% identical 2024-T3 plates, k_f/k_p = 0.70.
\begin{tikzpicture}[>=Latex,line join=round]
\tikzset{
  pa/.style={fill=black!10,draw=black!60,thick},
  pb/.style={fill=black!20,draw=black!60,thick},
  blt/.style={fill=blue!18,draw=blue!60!black,thick},
  ld/.style={->,black!75,line width=1.3pt},
  lb/.style={font=\scriptsize},
  sm/.style={font=\tiny,text=black!60}
}
\def\s{2.1}       % fastener spacing on the drawing

% ================= the splice =================
\begin{scope}[yshift=0cm]
\fill[pa] (-2.2,0.10) rectangle (9.2,0.70);
\draw[black!60,thick] (-2.2,0.10) rectangle (9.2,0.70);
\fill[pb] (-0.8,-0.70) rectangle (10.6,-0.10);
\draw[black!60,thick] (-0.8,-0.70) rectangle (10.6,-0.10);
\foreach \i in {0,1,2,3,4}{
  \pgfmathsetmacro\xx{\i*\s}
  \fill[blt] (\xx-0.17,-0.86) rectangle (\xx+0.17,0.86);
  \node[sm,anchor=south] at (\xx,0.95) {\the\numexpr\i+1\relax};
}
\draw[ld] (-3.1,0.40) -- (-2.3,0.40);
\node[lb,anchor=east] at (-3.2,0.40) {$P$};
\draw[ld] (10.7,-0.40) -- (11.5,-0.40);
\node[lb,anchor=west] at (11.6,-0.40) {$P$};
\node[sm,anchor=east] at (-2.3,0.95) {plate A};
\node[sm,anchor=west] at (10.7,-0.95) {plate B};
\end{scope}

% ================= the bar chart =================
\begin{scope}[yshift=-5.6cm]
\def\h{0.105}      % percent -> cm
\draw[->,black!70] (-0.9,0) -- (9.5,0);
\draw[->,black!70] (-0.9,0) -- (-0.9,4.6);
\node[lb,anchor=south,rotate=90] at (-2.05,2.3) {share of the load transferred};
\foreach \p in {0,10,20,30}{
  \draw[black!60] (-0.95,\p*\h) -- (-0.9,\p*\h);
  \node[sm,anchor=east] at (-0.98,\p*\h) {\p\,\%};}
% equal-share reference
\draw[dashed,black!55] (-0.9,20*\h) -- (9.3,20*\h);
\node[sm,anchor=west,fill=white,inner sep=1pt] at (6.25,20*\h+0.17) {equal share, $P/5$};
% bars
\foreach \i/\v in {0/34.2, 1/12.2, 2/7.2, 3/12.2, 4/34.2}{
  \pgfmathsetmacro\xx{\i*\s}
  \fill[blue!45!black,opacity=0.78] (\xx-0.46,0) rectangle (\xx+0.46,\v*\h);
  \draw[blue!45!black] (\xx-0.46,0) rectangle (\xx+0.46,\v*\h);
  \node[lb,anchor=south] at (\xx,\v*\h+0.04) {\v\,\%};
  \node[sm,anchor=north] at (\xx,-0.08) {\the\numexpr\i+1\relax};}
\end{scope}

% ================= the reading =================
\node[font=\footnotesize,align=left,text width=16.6cm,anchor=north west] at (-2.4,-6.35)
 {\textbf{Fasteners in a line do not share the load.} The outer pair here carry
  \SI{34}{\percent} each --- \textbf{1.7 times} an equal share --- while the middle one carries
  \SI{7}{\percent}. The reason is stiffness, not workmanship: at the first fastener, plate A still
  holds the whole load and plate B holds none, so the relative slip across the joint is greatest
  there. Two consequences you will use:
  \\[4pt]
  \textbf{1. Adding fasteners to a row buys less than you think.} Go to seven and the middle one
  takes \SI{2}{\percent}. Past about four, extra fasteners in the same line mostly add holes.
  \\[2pt]
  \textbf{2. The end fastener is the fatigue critical one.} It sees the highest bearing load
  \emph{and} the highest bypass stress, which is why cracks in a lap splice start at the outer
  rows --- exactly as the full-scale testing quoted in Lecture 5 found.};

\end{tikzpicture}
